% theme: familiar_substances
\MajorQuestion{身のまわりの物質}
\SetRowSpaceQanda

\Instruction{ガスバーナーと上皿てんびんの使い方に関する次の問いに答えなさい。}
\QQ{ガスバーナーで、炎の大きさを調節するねじを何というか。}{ガス調節ねじ}
\QQ{ガスバーナーで、炎に混じる空気の量を調節するねじを何というか。}{空気調節ねじ}
\QQ{ガスバーナーの炎は、空気の量を調節して何色の安定した炎にするか。}{青色}
\QQ{右ききの人が上皿てんびんでものの質量をはかるとき、はかるものは左右どちらの皿にのせるか。}{左の皿}
\QQ{上皿てんびんの分銅をのせたり取ったりするときに使う器具は何か。}{ピンセット}

\Instruction{電子てんびんとメスシリンダーの使い方、および物体と物質に関する次の問いに答えなさい。}
\QQ{電子てんびんでものの質量をはかる前に、表示を何\,gに合わせるか。}{0}%薬包紙を答えさせる問題に変更
\QQ{メスシリンダーの目盛りを読むとき、目の位置を何と同じ高さにするか。}{液面}
\QQ{メスシリンダーでは、最小目盛りの何分の1まで目分量で読み取るか。}{10分の1}
\QQ{ものの形や外観に注目したときの呼び名を\NamedBlank{object}という。}{物体}
\QQ{\RefBlank{object}の材料に注目したときの呼び名は何か。}{物質}

\Instruction{物質の分類について、次の問いに答えなさい。}
\QQ{金、銀、銅、鉄、アルミニウムなどの物質をまとめて\NamedBlank{metal}という。}{金属}
\QQ{ガラス、食塩、プラスチック、木、紙、ゴムなど、\RefBlank{metal}以外の物質を何というか。}{非金属}
\QQ{金属をみがいたときに現れる、特有のかがやきを何というか。}{金属光沢}
\QQ{物質が電気をよく通す性質を何というか。}{電気伝導性}%延性を答えさせる問題に変更
\QQ{物質が熱をよく伝える性質を何というか。}{熱伝導性}%展性を答えさせる問題に変更

\Instruction{物質の質量などに関する次の問いに答えなさい。}
\QQ{一定の体積あたりの質量を、その物質の\NamedBlank{density}という。}{密度}
\QQ{\RefBlank{density}を求める式を、質量と体積を用いて表しなさい。}{質量÷体積}%実際に密度を求める問題に変更
\QQ{液体より\RefBlank{density}が大きい物体は、その液体中でどうなるか。}{沈む}%密度・質量から体積を求める問題に変更
\QQ{炭素を含み、燃やすと二酸化炭素や水ができる物質を\NamedBlank{organic}という。}{有機物}
\QQ{\RefBlank{organic}以外の物質を何というか。}{無機物}
